
%  Diagramme de l'architecture U-Net (TikZ)

\begin{figure}[H]
\centering
\resizebox{\textwidth}{!}{%
\begin{tikzpicture}[
    font=\small,
    % --- styles des blocs ---
    attn/.style={
        draw, fill=yellow!40, diamond, 
        minimum width=0.4cm, minimum height=0.4cm,
        inner sep=0pt, line width=0.6pt
    },
    conv/.style={
        draw, fill=blue!20, rounded corners=3pt,
        minimum width=1.0cm, minimum height=2.2cm,
        line width=0.7pt
    },
    bottle/.style={
        draw, fill=purple!25, rounded corners=3pt,
        minimum width=1.0cm, minimum height=1.4cm,
        line width=0.7pt
    },
    pool/.style={
        draw, fill=red!15, rounded corners=2pt,
        minimum width=0.35cm, minimum height=2.2cm,
        line width=0.6pt
    },
    up/.style={
        draw, fill=green!15, rounded corners=2pt,
        minimum width=0.35cm, minimum height=2.2cm,
        line width=0.6pt
    },
    head/.style={
        draw, fill=orange!25, rounded corners=3pt,
        minimum width=0.7cm, minimum height=2.2cm,
        line width=0.7pt
    },
    lbl/.style={font=\scriptsize, align=center},
    % --- flèches ---
    arr/.style={-{Stealth[length=4pt]}, line width=0.8pt},
    skip/.style={-{Stealth[length=4pt]}, dashed, line width=0.8pt, color=gray!70},
]

% espacement
\def\dx{1.55}   % entre blocs conv successifs
\def\dxp{0.55}  % largeur pool/up

%encodeur
% enc0  (in=5, out=32)  h=2.2
\node[conv, minimum height=2.2cm] (e0) at (0,0) {};
\node[lbl, above=2pt of e0] {32};

% pool0
\node[pool, right=0.12cm of e0] (p0) {};

% enc1  (out=64)  h=1.8
\node[conv, minimum height=1.8cm, right=0.12cm of p0] (e1) {};
\node[lbl, above=2pt of e1] {64};

% pool1
\node[pool, minimum height=1.8cm, right=0.12cm of e1] (p1) {};

% enc2  (out=128)  h=1.4
\node[conv, minimum height=1.4cm, right=0.12cm of p1] (e2) {};
\node[lbl, above=2pt of e2] {128};

% pool2
\node[pool, minimum height=1.4cm, right=0.12cm of e2] (p2) {};

% enc3  (out=256)  h=1.0
\node[conv, minimum height=1.0cm, right=0.12cm of p2] (e3) {};
\node[lbl, above=2pt of e3] {256};

% pool3
\node[pool, minimum height=1.0cm, right=0.12cm of e3] (p3) {};

%bottleneck
\node[bottle, right=0.12cm of p3] (bn) {};
\node[lbl, above=2pt of bn] {512};
\node[lbl, below=2pt of bn] {Dropout\\(p=0.3)};

% decodeur
% up3
\node[up, minimum height=1.0cm, right=0.12cm of bn] (u3) {};
% dec3  (out=256)  h=1.0
\node[conv, minimum height=1.0cm, right=0.12cm of u3] (d3) {};
\node[lbl, above=2pt of d3] {256};

% up2
\node[up, minimum height=1.4cm, right=0.12cm of d3] (u2) {};
% dec2  (out=128)  h=1.4
\node[conv, minimum height=1.4cm, right=0.12cm of u2] (d2) {};
\node[lbl, above=2pt of d2] {128};

% up1
\node[up, minimum height=1.8cm, right=0.12cm of d2] (u1) {};
% dec1  (out=64)  h=1.8
\node[conv, minimum height=1.8cm, right=0.12cm of u1] (d1) {};
\node[lbl, above=2pt of d1] {64};

% up0
\node[up, minimum height=2.2cm, right=0.12cm of d1] (u0) {};
% dec0  (out=32)  h=2.2
\node[conv, minimum height=2.2cm, right=0.12cm of u0] (d0) {};
\node[lbl, above=2pt of d0] {32};

% (Conv 1×1 → logits)
\node[head, right=0.3cm of d0] (hd) {};
\node[lbl, above=2pt of hd] {Conv\\$1{\times}1$};

% entrée et sortie
\node[lbl, left=0.35cm of e0] (inp) {Entrée\\$320{\times}320$\\5 bandes};
\draw[arr] (inp) -- (e0);

\node[lbl, right=0.35cm of hd] (out) {Masque\\$320{\times}320$\\(logits)};
\draw[arr] (hd) -- (out);

% flèches
\foreach \a/\b in {e0/p0, p0/e1, e1/p1, p1/e2, e2/p2, p2/e3, e3/p3,
                   p3/bn, bn/u3, u3/d3, d3/u2, u2/d2, d2/u1, u1/d1,
                   d1/u0, u0/d0, d0/hd}
    \draw[arr] (\a) -- (\b);


%skip connections
% Calcul des positions y pour les arcs
% enc0 → dec0  (niveau 0, h=2.2)
% enc3 → dec3 (avec Attention Gate)
\node[attn] (ag3) at ($(d3.north)+(0,0.55)$) {};
\draw[skip] (e3.north) to[out=90, in=90, looseness=0.40] (ag3.north);
\draw[arr, color=gray!70] (ag3) -- (d3.north);

% enc2 → dec2 (avec Attention Gate)
\node[attn] (ag2) at ($(d2.north)+(0,0.55)$) {};
\draw[skip] (e2.north) to[out=90, in=90, looseness=0.34] (ag2.north);
\draw[arr, color=gray!70] (ag2) -- (d2.north);

% enc1 → dec1 (avec Attention Gate)
\node[attn] (ag1) at ($(d1.north)+(0,0.55)$) {};
\draw[skip] (e1.north) to[out=90, in=90, looseness=0.30] (ag1.north);
\draw[arr, color=gray!70] (ag1) -- (d1.north);

% enc0 → dec0 (avec Attention Gate)
\node[attn] (ag0) at ($(d0.north)+(0,0.55)$) {};
\draw[skip] (e0.north) to[out=90, in=90, looseness=0.28] (ag0.north);
\draw[arr, color=gray!70] (ag0) -- (d0.north);


%légende
\begin{scope}[shift={(0,-2.2)}]
    \node[conv,   minimum height=0.45cm, minimum width=0.75cm] (le) at (0,0) {};
    \node[lbl, right=0.1cm of le] {DoubleConv (Conv $3{\times}3$ $\to$ BN $\to$ ReLU) $\times 2$};

    \node[pool,   minimum height=0.45cm, minimum width=0.35cm] (lp) at (0,-0.65) {};
    \node[lbl, right=0.1cm of lp] {MaxPool $2{\times}2$};

    \node[up,     minimum height=0.45cm, minimum width=0.35cm] (lu) at (3.5,-0.65) {};
    \node[lbl, right=0.1cm of lu] {ConvTranspose $2{\times}2$};

    \node[bottle, minimum height=0.45cm, minimum width=0.75cm] (lb) at (12,0) {};
    \node[lbl, right=0.1cm of lb] {Bottleneck + Dropout};

    % skip legend
    \draw[skip] (7.2,0) -- (8.0,0);
    \node[lbl, right=0.05cm] at (8.0,0) {Skip connection};
    \node[attn] (la) at (7.2,-0.65) {};
    \node[lbl, right=0.1cm of la] {Attention Gate};
\end{scope}

% labels encodeur/decodeur/bottleneck
\node[lbl, gray, above=0.55cm] at ($(e0)!0.5!(p3)$) {\textsc{Encodeur}};
\node[lbl, gray, above=0.55cm] at ($(u3)!0.5!(d0)$) {\textsc{Décodeur}};
\node[lbl, gray, above=0.55cm] at (bn)              {\textsc{Bottleneck}};


\foreach \n in {p0,p1,p2,p3}
    \node[lbl, rotate=90, font=\tiny, text=red!60!black] at (\n) {MP};
\foreach \n in {u0,u1,u2,u3}
    \node[lbl, rotate=90, font=\tiny, text=green!50!black] at (\n) {CT};

\end{tikzpicture}
}
\caption{Architecture U-Net avec Attention Gates utilisée pour la segmentation glaciaire.}
\label{fig:unet_architecture}
\end{figure}